\subsection{CX-N15: a complete accessible composition for general factorization}
\label{finding:CX-N15}

\paragraph{Finding and source interface.}
Assume \(\mathsf H\), the all-number-field Hecke claim.
Given an explicit finite field \(\mathbb F_q\), \(q=p^m\), and
\(0\ne f\in\mathbb F_q[X]\) of degree \(d\geq2\), there is a
deterministic algorithm that returns its complete factorization in
\[
 \bigl(d^{\log d}\log q\bigr)^{O(1)}
\]
bit operations.
The composition below uses Adleman--Lenstra Algorithm~A and
Theorem~1, with its auxiliary-prime step supplied by CX-N11;
Bach--von zur Gathen--Lenstra Theorem~2; and the IKRS Main Theorem,
together with the explicitly proved coefficient-unit step of
CX-N14
\cite{CXNAdlemanLenstra86,CXNBGL01,CXNIKRS08}.
It does not require the unretrieved internal reduction of
Evdokimov's 1994 paper, or a general finite-field factorization
oracle.
This gives an independent proof of the broad transfer in CX-N09;
it does not change that theorem's degree dependence.

\paragraph{Step 1: construct a maximal coefficient unit.}
Consider a split algebra \(A_0\cong\mathbb F_p^{d_0}\) with a
supplied nonidentity automorphism and \(p>d\geq d_0\).
Reduce its automorphism to prime order \(r\leq d_0\) as in
CX-N14. The branch \(p\leq d\) is handled unconditionally in
Step~3. The following construction works for every such \(r\),
including \(r\mid p-1\). Put
\[
 s=\operatorname{ord}_r(p),\qquad Q=p^s,
 \qquad e=v_r(Q-1),\qquad R=r^e.
\]
The quantities \(s\leq r-1\), \(e\), and \(R\) are computable
in polynomial time in \(r,\log p\); no factorization of the
complementary part of \(Q-1\) is used.
Adleman--Lenstra Algorithm~A, with CX-N11, constructs an explicit
field
\(U=\mathbb F_{p^{sr}}\) in polynomial time in \(r,\log p\).
Indeed \(sr\leq r(r-1)\) and \(p\nmid sr\), because both
\(r\) and \(s\) are smaller than \(p\).
Compute the subfield
\[
 B=\ker\bigl(U\longrightarrow U,\ u\longmapsto u^Q-u\bigr)
   \cong\mathbb F_Q
\]
by \(\mathbb F_p\)-linear algebra, and compute its multiplication
table and inclusion in \(U\).
This is exactly the explicit-data convention of BGL, adopted
from Lenstra: basis structure constants suffice, and Frobenius
kernels provide explicit subfields
\cite[p.~329 and \S2, p.~332]{CXNLenstra91}.

BGL Theorem~2, p.~7, constructs a primitive \(r\)-th root
\(\zeta\in B\), given additionally an irreducible polynomial
over \(\mathbb F_p\) of degree \(\ell\) for each prime
\(\ell\mid r-1\) with \(\ell\nmid s\).
All these primes are smaller than \(p\) and at most \(r-1\).
Trial division factors \(r-1\) in polynomial time in \(r\),
and Adleman--Lenstra with CX-N11 constructs every required
auxiliary polynomial in polynomial time in \(r,\log p\).
Thus all BGL inputs are supplied, and its output and running-time
conditions apply. When \(r=2\), simply take \(\zeta=-1\);
there are no auxiliary primes.

Find a nonzero solution of
\[
 u^Q=\zeta u\qquad (u\in U)
\]
by \(\mathbb F_p\)-linear algebra.
Such a solution exists: the image of the power map
\(U^\times\to U^\times\), \(v\mapsto v^{Q-1}\), is the
cyclic subgroup of order
\[
 S=\frac{Q^r-1}{Q-1}=1+Q+\cdots+Q^{r-1}.
\]
Since \(Q\equiv1\pmod r\), the integer \(S\) is divisible
by \(r\); hence this subgroup contains \(\zeta\).
Let \(c=u^r\). Then
\[
 c^Q=(\zeta u)^r=c,
 \qquad c^{(Q-1)/r}=u^{Q-1}=\zeta\ne1.
\]
Thus \(c\in B\) is an \(r\)-th nonresidue.
Because \(B^\times\) is cyclic of order \(Q-1=R w_0\),
\(r\nmid w_0\), the element
\[
 \eta_0=c^{w_0}\in B^\times
\]
has exact order \(R\).
To verify the last assertion, write \(c=g^a\) for an abstract
generator \(g\) of \(B^\times\).
The nonresidue test says \(r\nmid a\), so
\(\operatorname{ord}(g^{a w_0})=R\).
This argument uses a generator only to prove the order; none is
an algorithmic input.

\paragraph{Step 2: descend the unit by equivariant interpolation.}
Let
\[
 D_r=\mathbb F_p[T]/(\Phi_r(T)),\qquad
 G=(\mathbb Z/r\mathbb Z)^\times,
 \qquad \rho_b(T)=T^b\quad(b\in G),
\]
and put \(\omega(a)=a^{r^{e-1}}\bmod R\).
This is the multiplicative Teichm\"uller lift of \(G\); its
values have order dividing \(r-1\) and reduce to \(a\) modulo
\(r\).
Since \(p^s\equiv1\pmod R\) and \(s\mid r-1\),
\(p\bmod R\) has order \(s\) and equals
\(\omega(p\bmod r)\).
Indeed reduction modulo \(r\) is injective on the
prime-to-\(r\) torsion of \((\mathbb Z/R\mathbb Z)^\times\).

Interpolate the unique polynomial \(P\in B[T]\),
\(\deg P<r-1\), with values
\[
 P(\zeta^a)=\eta_0^{\omega(a)}\qquad(a\in G).
\]
There are only \(r-1\) distinct nodes, so interpolation has
polynomial cost.
Write \(P^{(p)}\) for the polynomial obtained by raising each
coefficient to the \(p\)-th power.
For every \(a\in G\), the Frobenius identity and the preceding
Teichm\"uller congruence give
\[
 P^{(p)}(\zeta^{pa})
   =P(\zeta^a)^p
   =\eta_0^{p\omega(a)}
   =\eta_0^{\omega(pa)}
   =P(\zeta^{pa}).
\]
Both polynomials have degree below \(r-1\) and agree on all
\(r-1\) nodes. Hence \(P^{(p)}=P\), and
\(P\in\mathbb F_p[T]\).
These coefficients are extracted on the known prime-field
subspace of \(B\).

Set \(\eta=P(T)\bmod\Phi_r(T)\in D_r\).
Evaluation at all roots \(\zeta^a\) is injective after extending
coefficients to \(B\), because \(\Phi_r\) is separable and
splits there.
Every evaluated value is a unit of order \(R\), so
\(\eta^R=1\) and \(\operatorname{ord}(\eta)=R\), the latter
already following from evaluation at \(\zeta\).
For \(a,b\in G\),
\[
 P(\zeta^{ab})
 =\eta_0^{\omega(a)\omega(b)}
 =P(\zeta^a)^{\omega(b)}.
\]
Injectivity of evaluation therefore proves
\(\rho_b(\eta)=\eta^{\omega(b)}\).
Thus \(\eta\) is the maximum-order coefficient Teichm\"uller
unit required by CX-N14, constructed in
\(\operatorname{poly}(d_0,\log p)\) time under \(\mathsf H\).
No factorization of \(\Phi_r\) or general polynomial-factorization
oracle has been used.

For completeness, the last step of CX-N14 applies here for
every \(r\ne p\), without its prime-field nonresidue shortcut.
Its moved-component ideal either already splits \(A_0\), or
establishes a semiregular prime-order automorphism \(\tau\).
IKRS Proposition~4.6 then supplies
\(x\in\mathcal T_{A_0,r}\) with
\(\tau(x)=z x\), where \(z=T\bmod\Phi_r(T)\).
The algebra \(A_0\otimes D_r\) is a product of copies of
\(\mathbb F_Q\), so all of its units of \(r\)-power order have
order dividing \(R\).
The coefficient unit \(\eta\) has order \(R\) and is fixed
by \(\tau\). If \(\langle x,\eta\rangle\) were cyclic,
\(\eta\) would generate it, contradicting
\(\tau(x)=zx\ne x\).
IKRS Lemma~4.2 therefore returns a zero divisor in \(A_0\)
in polynomial time.
This proves the claimed splitting step for all prime orders in
the branch \(p>d\), including \(s=1\) and \(r=2\), without
adding an unstated Dirichlet-character premise.

\paragraph{Step 3: from a general polynomial to a split algebra.}
Polynomial gcds and derivatives give squarefree preprocessing.
If a positive-degree polynomial has zero derivative, it is a
\(p\)-th power: take the coefficient \(p\)-th roots using the
inverse Frobenius \(a\mapsto a^{p^{m-1}}\), divide the exponents
by \(p\), and retain the multiplicity \(p\).
Otherwise a nontrivial gcd with the derivative yields smaller
factors to process recursively.
These operations have polynomial bit complexity.
It therefore suffices to factor a monic squarefree polynomial
\(f\) of degree at most \(d\).

Compute the multiplication table of
\(A=\mathbb F_q[X]/(f)\), viewed as an
\(\mathbb F_p\)-algebra of dimension \(m\deg f\), and then
compute
\[
 B_f=\ker\bigl(A\longrightarrow A,\ a\longmapsto a^p-a\bigr).
\]
The map is \(\mathbb F_p\)-linear, so its kernel and the kernel's
multiplication table are obtained by linear algebra.
Abstractly, if \(f=\prod_{i=1}^t f_i\) is its distinct
irreducible factorization, then
\[
 A\cong\prod_{i=1}^t\mathbb F_{q^{\deg f_i}},
 \qquad B_f\cong\mathbb F_p^t.
\]
This proves the split promise without knowing the factors.
The primitive idempotents of \(B_f\) and \(A\) coincide.
If \(\dim_{\mathbb F_p}B_f=1\), the polynomial is irreducible.

If \(p\leq d\), choose \(b\in B_f\setminus\mathbb F_p1\)
and enumerate \(b-a1\) for \(a\in\mathbb F_p\).
Each is nonzero, and at least one has a zero component in
\(B_f\cong\mathbb F_p^t\).
Testing multiplication matrices therefore finds a nonzero zero
divisor after at most \(p\leq d\) tests.
This early branch works even when \(t<p\).

For \(p>d\), apply the IKRS Main Theorem, p.~3, to the explicit
\(t\)-dimensional algebra \(B_f\).
It deterministically gives a decomposition into ideals, each
with an automorphism whose order equals that ideal's dimension,
in time \(\operatorname{poly}(t^{\log t},\log p)\).
If the decomposition has more than one summand, the identity of
any proper nonzero summand is a nonzero zero divisor of \(B_f\);
that identity is found by linear algebra in the supplied ideal.
If there is only one summand, its automorphism has order
\(t>1\), and the preceding steps with CX-N14 produce a nonzero
zero divisor of \(B_f\) in polynomial additional time.
Thus in either case we obtain a nonzero zero divisor \(b\in A\),
and \(\gcd(f,b)\), using its polynomial representative, is a
proper nonconstant factor.
Equivalently, a primitive idempotent \(e_i\) identifies the
component factor as \(\gcd(f,1-e_i)\).
Recurse until all factors are irreducible, and reinsert the
recorded multiplicities.

\paragraph{Running time and exact analytic scope.}
There are at most \(2d-1\) nodes in the nontrivial factorization
tree, and every split algebra has dimension at most \(d\).
Linear algebra in the original quotient has size polynomial in
\(d,m,\log p\), hence in \(d,\log q\).
At each node the coefficient construction has polynomial cost;
the IKRS call has cost
\(d^{O(\log d)}(\log q)^{O(1)}\).
Multiplying by the number of nodes preserves the displayed
\((d^{\log d}\log q)^{O(1)}\) bound.

The analytic assumptions enter the constructed fields and
auxiliary irreducible polynomials through Adleman--Lenstra's
auxiliary-prime step only.
An explicit sufficient family is RH for \(\zeta_{\mathbb Q}\)
and for
\[
 \zeta_{E_\nu},\qquad
 E_\nu=\mathbb Q(\zeta_{2\nu},p^{1/\nu}),
\]
where \(\nu=sr\) or \(\nu=\ell\) for one of the auxiliary
primes \(\ell\mid r-1\), \(\ell\nmid s\), used at a node.
CX-N11 supplies the required auxiliary prime from precisely
these Dedekind-zeta hypotheses, all implied by \(\mathsf H\).
Every \(\nu\leq d^2\) is coprime to \(p\), because these
constructions occur only after the branch \(p\leq d\) has
already been handled, and the factors of \(sr\) are smaller
than \(p\).
This is a sufficient family, not a necessity assertion.
BGL's root-of-unity algorithm with those inputs, the
Frobenius/eigenvector calculations, interpolation, and IKRS
are unconditional algebraic steps.
The construction therefore supplies the formerly isolated
coefficient-unit gate under \(\mathsf H\) without circularly
calling general factorization or assuming a polynomial bound
for its varying degree.
\hfill\(\square\)

\paragraph{Finite-\'etale restatement.}
For squarefree \(f\), \(\operatorname{Spec}A\) is a finite
\'etale scheme over \(\mathbb F_p\).
Its Frobenius-fixed function algebra \(B_f\) is the constant
finite \'etale algebra indexing its connected components.
Idempotents found in \(B_f\) pull back to the corresponding
open-and-closed decomposition of \(\operatorname{Spec}A\).
The proof constructs these component projectors with the stated
quasipolynomial dependence on the input polynomial degree.
